% =====================================================================
\chapter{Готовые рецепты}
% =====================================================================
\chapterintro
  {полные настройки для типовых моделей: FPV-квадрокоптера на Betaflight, учебного
   самолёта, летающего крыла, планера с кроу-тормозом, машины или лодки}
  {всё из предыдущих глав}

\section{FPV-квадрокоптер (Betaflight + ELRS)}

\subsection{Модель в пульте}

\begin{enumerate}
  \item Создайте модель, имя \code{Quad-5in}. \menu{Internal RF} = \menu{CRSF}, привяжите
        приёмник (глава~\ref{ch:rf}). Packet Rate 250\,Гц, Telem Ratio \menu{Std},
        Switch Mode \menu{Wide}.
  \item Входы \sw{I1}--\sw{I4}: вес 100, экспо 0, \menu{Trim} = \menu{OFF}.
  \item Миксеры по таблице ниже (порядок AETR).
  \item Таймеры: Timer1 --- обратный отсчёт 4:00 по арму, Timer2 --- суммарный налёт.
  \item Предполётные проверки: газ внизу, \sw{SA} вверху (дизарм).
  \item Телеметрия: \menu{Discover new sensors}, экран \menu{Nums}, предупреждение о батарее.
\end{enumerate}

\begin{center}\small
\renewcommand{\arraystretch}{1.15}
\begin{tabular}{@{}llll@{}}
\toprule
Канал & TX12 & Boxer & Функция в Betaflight \\
\midrule
\sw{CH1}--\sw{CH4} & стики & стики & Roll, Pitch, Throttle, Yaw \\
\sw{CH5} (AUX1) & \sw{SA} & \sw{SA} & \textbf{ARM} (только 2 положения!) \\
\sw{CH6} (AUX2) & \sw{SB} & \sw{SB} & ANGLE / HORIZON / (acro) \\
\sw{CH7} (AUX3) & \sw{SC} & \sw{SC} & BEEPER (найти потерянный коптер) \\
\sw{CH8} (AUX4) & \sw{SE} & \sw{SD} & FLIP OVER AFTER CRASH («черепаха») \\
\sw{CH9} (AUX5) & \sw{SF} & \sw{SE} & запись видео / PREARM / LAUNCH CONTROL \\
\bottomrule
\end{tabular}
\end{center}

\subsection{В Betaflight Configurator}

\begin{enumerate}
  \item Вкладка \menu{Receiver}: \menu{Serial (via UART)}, провайдер \menu{CRSF};
        \menu{Channel map} = \code{AETR1234}. Пошевелите стиками --- полоски должны двигаться
        правильно, в центре --- около 1500, по краям --- около 1000 и 2000.
  \item Вкладка \menu{Modes}: \menu{ARM} на AUX1 в диапазоне 1700--2100 (верхнее или нижнее
        положение --- как вам удобно), \menu{ANGLE} на AUX2 и т.\,д.
  \item Вкладка \menu{Failsafe}: проверьте, что при выключении пульта коптер переходит в
        failsafe (со снятыми пропеллерами!).
  \item Для голосовых сообщений о напряжении на банку включите в CLI
        \code{set report\_cell\_voltage = ON}.
\end{enumerate}

\begin{recipe}{озвучка квадрокоптера}
\begin{itemize}
  \item \sw{SA↓} → \menu{Play track} \code{armed}; \sw{SA↑} → \code{disarm}.
  \item \sw{SB} каждое положение → \code{angle}, \code{horizon}, \code{acro}.
  \item \sw{L1}: \code{a<x} \sw{RxBt} 3,5\,В, \menu{Delay} 2\,с → \menu{Play track}
        \code{batlow}, повтор 10\,с.
  \item \sw{L2}: \code{a<x} \sw{RQly} 70 → \menu{Haptic} + \menu{Play value} \sw{RQly}, повтор 5\,с.
\end{itemize}
\end{recipe}

\section{Учебный самолёт (4 канала, PWM-приёмник)}

\begin{center}\small
\begin{tabular}{@{}lll@{}}
\toprule
Канал & Источник & Настройки \\
\midrule
\sw{CH1} & \sw{I1} Ail & расходы на \sw{SC}: 100/70/50\,\%, экспо 25--35\,\% \\
\sw{CH2} & \sw{I2} Ele & расходы на \sw{SC}: 100/70/50\,\%, экспо 25--35\,\% \\
\sw{CH3} & \sw{I3} Thr & отсечка газа \menu{Override} на \sw{SD} (TX12) / \sw{SA} (Boxer) \\
\sw{CH4} & \sw{I4} Rud & экспо 20\,\% \\
\bottomrule
\end{tabular}
\end{center}

\begin{itemize}
  \item Настройте в \menu{Outputs} направления и пределы (глава 9). Проверьте:
        ручка вправо --- правый элерон вверх; ручка на себя --- руль высоты вверх;
        правая педаль --- руль направления вправо.
  \item Failsafe для PWM-приёмника: газ --- минимум, рули --- нейтраль или лёгкий
        разворот.
  \item Таймер \menu{Throttle \%} на 8--10 минут: учитывает реальный расход батареи.
\end{itemize}

\section{Летающее крыло}

\begin{itemize}
  \item Миксеры элевонов --- глава~\ref{ch:mixes}. Веса 50/50; для ручного полёта крылу
        обычно нужно больше тангажа, чем крена: Ele 60, Ail 40.
  \item Экспо 30--40\,\% --- у крыла очень резкая реакция у центра.
  \item Если стоит полётный контроллер INAV/ArduPilot, \emph{не} делайте микширование в пульте:
        передавайте чистые Ail/Ele, а смешивает контроллер.
  \item Отдельный «запуск»: \sw{SE} (или кнопка) → полётный режим \sw{FM1} с небольшим
        «кабрированием» (Ele $+10\,\%$) и таймером на запуск мотора.
\end{itemize}

\section{Планер с кроу-тормозом}

Планер с двумя элеронами (\sw{CH1}, \sw{CH6}) и двумя закрылками (\sw{CH5}, \sw{CH7}),
мотор на \sw{CH3} управляется тумблером или стиком. Кроу (\emph{crow}, «бабочка») ---
элероны вверх, закрылки вниз --- резко гасит скорость на посадке.

\begin{center}\small
\renewcommand{\arraystretch}{1.15}
\begin{tabular}{@{}lllll@{}}
\toprule
Канал & Source & Weight & Curve & Switch / FM \\
\midrule
\sw{CH1} AilL & \sw{I1} Ail & 100 & Diff 30 & \\
              & \sw{Thr} (стик) & 50 & x>0 & FM Land \\
\sw{CH6} AilR & \sw{I1} Ail & $-100$ & Diff 30 & \\
              & \sw{Thr} (стик) & 50 & x>0 & FM Land \\
\sw{CH5} FlpL & \sw{Thr} (стик) & $-80$ & x>0 & FM Land \\
\sw{CH7} FlpR & \sw{Thr} (стик) & $-80$ & x>0 & FM Land \\
\sw{CH2} Ele  & \sw{I2} Ele & 100 & & \\
              & \sw{Thr} (стик) & $-10$ & x>0 & FM Land (компенсация) \\
\bottomrule
\end{tabular}
\end{center}

В режиме «Land» стик газа управляет кроу, а мотор принудительно выключен
(\menu{Override CH3} $-100$ в этом режиме). Знаки весов зависят от установки серв ---
проверяйте на земле.

\section{Машина или лодка}

\begin{itemize}
  \item \sw{CH1} --- руль (правый стик, горизонталь), \sw{CH2} --- газ (левый стик,
        вертикаль; у многих водителей газ с центровкой --- вперёд/назад). Если у вас пульт
        с трещоткой на газу, её можно ослабить или заменить пружиной (у Boxer и TX12
        это настраивается винтами стика).
  \item Ограничитель скорости для ребёнка или новичка: \sw{GV1} = 50 и вес газа в миксере =
        \code{GV1}; \menu{Adjust GV1} от \sw{S1} на переключателе --- «родительский контроль».
  \item Экспо на руле 20--30\,\%, расходы руля на переключателе.
  \item Failsafe: газ в нейтраль (не в минимум! у машины с реверсом минимум --- это полный
        назад).
\end{itemize}

\begin{summary}
Для любой модели порядок один: радиомодуль и привязка → входы → миксеры → выходы
(направления и пределы) → таймеры и предполётные проверки → failsafe → телеметрия и голос →
проверка на земле → первый полёт.
\end{summary}
